% concatenated from top-FLM.tex
\documentclass[11pt]{article}
\usepackage[T1]{fontenc}
\usepackage{lmodern}
\usepackage[margin=1in]{geometry}
\usepackage{amsmath,amssymb,amsthm,mathtools}
\usepackage{microtype}

\usepackage[hidelinks]{hyperref}
\hypersetup{pdftitle={Superpolynomial lower bounds for vertex numbers of real projective space triangulations and a topological Figiel--Lindenstrauss--Milman theorem}}
\usepackage{xcolor}

\newtheorem{theorem}{Theorem}
\newtheorem{corollary}[theorem]{Corollary}
\newtheorem{lemma}[theorem]{Lemma}
\theoremstyle{remark}
\newtheorem{remark}[theorem]{Remark}
\DeclareMathOperator{\ind}{ind}
\DeclareMathOperator{\coind}{coind}
\DeclareMathOperator{\Hess}{Hess}
\DeclareMathOperator{\tr}{tr}
\DeclareMathOperator{\diag}{diag}
\newcommand{\R}{\mathbb R}
\newcommand{\Z}{\mathbb Z}
\newcommand{\RP}{\mathbb {RP}}
\newcommand{\Czero}{C_0}
\newcommand{\nuplus}{\nu_{\geq 0}}

\title{Superpolynomial lower bounds for vertex numbers of\\
real projective space triangulations via a\\
topological Figiel--Lindenstrauss--Milman theorem}

\author{
Florian Frick\thanks{Department of Mathematical Sciences, Carnegie Mellon University. Email: \texttt{frick@cmu.edu}.}
\and
Kaave Hosseini\thanks{Department of Computer Science, University of Rochester. Email: \texttt{kaave.hosseini@rochester.edu}.}
\and
Eric Myzelev\thanks{Department of Mathematical Sciences, Carnegie Mellon University. Email: \texttt{etmyzele@andrew.cmu.edu}.}
\and
Arya Narnapatti\thanks{Department of Mathematical Sciences, Carnegie Mellon University. Email: \texttt{anarnapa@andrew.cmu.edu}.}
\and
Aliaksei Vasileuski\thanks{Department of Mathematical Sciences, Carnegie Mellon University. Email: \texttt{avasileu@andrew.cmu.edu}.}
}

\date{September 9, 2026}

\begin{document}
\maketitle
\vspace{-1.5em}

\begin{abstract}
We prove that every simplicial triangulation of real projective
$d$-space has $\exp(\Omega(\sqrt d))$ vertices. Together with known
constructions, this determines the minimum vertex number as
$\mu_d=\exp(d^{1/2+o(1)})$. The result follows from a topological generalization of the
Figiel--Lindenstrauss--Milman inequality, answering a recent question of Frick,
Hosseini, and Vasileuski: a finite strongly regular CW
complex with a free cellular involution, $v$ vertices, and $f$ maximal
cells has $\Z/2$-index at most $O(\log v\log f)$. We bound the dimensions
of Morse cells by a trace estimate for a constrained Hessian, obtaining
a Morse-theoretic proof of the classical inequality for centrally
symmetric polytopes. As further applications of this
inequality, we give an $\exp(\Omega(\sqrt t))$ lower bound for the order
of a triangle-free topologically $t$-chromatic graph and bound the index of sign complexes by $O(d\log^2 N)$ for total
matrices and $O(d\log^3 N)$ for partial matrices, where $N\geq2$
is the number of columns and $d\geq1$ is the VC dimension.
\end{abstract}

\section{Introduction}

How few vertices can a simplicial triangulation of $\RP^d$ have?
Write $\mu_d$ for this minimum. The classical lower bound of
Arnoux--Marin~\cite{AM} is $\mu_d\geq\binom{d+2}{2}+1$ for $d\geq3$.
The subexponential constructions of
Adiprasito--Avvakumov--Karasev~\cite{AAK}, with the improvement of
Frankl--Pach--P\'alv\"olgyi~\cite[equation~(3)]{FPP}, give
$\mu_d\leq\exp(O(\sqrt{d\log d}))$.
We prove the superpolynomial lower bound
$\mu_d\geq\exp(c\sqrt d)$ for an absolute $c>0$ and all sufficiently
large $d$. Consequently,
\[
 c\sqrt d\leq\log\mu_d\leq C\sqrt{d\log d},
 \qquad \mu_d=\exp(d^{1/2+o(1)})
 \qquad(d\longrightarrow\infty).
\]
The power $1/2$ of $d$ in the exponent is thus determined; a factor
$\sqrt{\log d}$ remains between the bounds on $\log\mu_d$.

Our proof uses a topological extension of the
Figiel--Lindenstrauss--Milman inequality. The classical inequality
states that a centrally symmetric $d$-dimensional polytope with
$v$ vertices and $f$ facets satisfies
\begin{equation}\label{eq:flm}
 \log v\,\log f\geq c d
\end{equation}
for an absolute constant $c>0$ \cite{FLM}.
Frick, Hosseini, and Vasileuski \cite[Question~37]{FHV} asked
whether an analogous bound holds for finite strongly regular
CW complexes with a free cellular involution, with equivariant
index replacing dimension.

A finite CW complex is \emph{strongly regular} if its characteristic
maps are embeddings and the intersection of two closed cells is
empty or a closed cell. Its \emph{facets} are its inclusion-maximal
cells. We answer the question affirmatively. The proof combines
smooth Morse theory with an elementary trace estimate and also
gives a Morse-theoretic proof of \eqref{eq:flm}. For geometric
proofs and further context, see \cite{FLM,Milo}.

All logarithms are natural. For a space $X$ with a free involution~$\tau$ (that is, $\tau(x)\ne x$ for all~$x$), 
write
\[
 \ind X=\min\{k\geq0:\text{there is an equivariant map }X\longrightarrow S^k\},
\]
where spheres carry the antipodal involution; the value is $+\infty$
if no such map exists, and $\ind\varnothing=-1$. Equivariant here means that the map commutes with free involutions on the domain and codomain.
The index is monotone under equivariant maps, and $\ind S^d=d$ by the
Borsuk--Ulam theorem; see Matou\v sek's book~\cite{Matousek} for an introduction.
The corresponding \emph{coindex} is
\[
 \coind X=\sup\{k\geq0:\text{there is an equivariant map }
 S^k\longrightarrow X\}.
\]
We set $\coind\varnothing=-1$ and allow the value $+\infty$.
The Borsuk--Ulam theorem gives $\coind X\leq\ind X$.

\begin{theorem}[Topological FLM inequality]\label{thm:main}
Let $X$ be a nonempty finite strongly regular CW complex with a free
involution that permutes its cells. If $X$ has $v$ vertices and $f$
facets, then
\begin{equation}\label{eq:main}
 \ind X\leq \Czero\log v\,\log f,
 \qquad \Czero=24e^2.
\end{equation}
In particular, this holds for every finite simplicial complex with a
free simplicial involution.
\end{theorem}

If $X=\partial P$ for a centrally symmetric
$d$-polytope, radial projection gives an equivariant homeomorphism
$X\cong S^{d-1}$. Hence $\ind X=d-1$, and \eqref{eq:main} gives
$\log v\,\log f\geq(d-1)/C_0$.
Thus \eqref{eq:flm} holds with $c=(2C_0)^{-1}$.
\paragraph{Applications.}
To apply Theorem~\ref{thm:main} to a triangulation of $\RP^d$
with $N$ vertices, we lift it to an antipodal triangulation of
$S^d$ and enlarge the lift to a free simplicial complex with
$2N$ vertices and at most $2N$ facets. This gives
$d\leq\Czero\log^2(2N)$ and hence the following bound; the
construction appears in Section~\ref{subsection:triangulations}.
\begin{corollary}[Projective-space triangulations]\label{cor:rp}
Every finite simplicial triangulation of $\RP^d$, $d\geq1$, has at least
\begin{equation}\label{eq:rp}
 N\geq \frac12\exp\!\left(\sqrt{d/\Czero}\right)
\end{equation}
vertices.
\end{corollary}
Frankl--Pach--P\'alv\"olgyi \cite[Theorem~1]{FPP} already proved an
$\exp(\Omega(\sqrt d))$ lower bound for the sizes of the exchange
set systems underlying the construction of Adiprasito--Avvakumov--Karasev \cite{AAK}.
Corollary~\ref{cor:rp} applies more generally to every simplicial triangulation of~$\RP^d$.

It also answers \cite[Question~3.8]{AAK}, concerning centrally symmetric
polytopes whose facets through $v$ are disjoint from those through $-v$
for every vertex $v$. Such a polytope admits an equivariant boundary
triangulation whose antipodal quotient is simplicial, without adding
vertices \cite[Corollary~2.2]{AAK}. Its number of vertices is therefore
also $\exp(\Omega(\sqrt d))$ in dimension $d$.

More generally, the same vertex bound holds for any finite simplicial
complex $T$ admitting a double cover of Stiefel--Whitney height at
least $d$, equivalently a class $\alpha\in H^1(T;\mathbb F_2)$ with
$\alpha^d\ne0$; see Corollary~\ref{cor:cup}.

For a graph $G$, let $B_0(G)$ be its box complex, with empty shores
allowed; its definition is recalled in Section~\ref{sec:graphs}.
For an integer $t\geq1$, following Simonyi--Tardos
\cite[Definition~3]{ST}, $G$ is \emph{topologically
$t$-chromatic} if there is an equivariant map $S^{t-1}\to |B_0(G)|$.
This implies $\chi(G)\geq t$.

\begin{corollary}[Triangle-free graphs]\label{cor:graph}
If $G$ is a finite simple triangle-free graph on $n\geq1$ vertices, then
\begin{equation}\label{eq:graphindex}
 \ind |B_0(G)|\leq\Czero\log(2n)\log(2n+2).
\end{equation}
In particular, if $G$ is topologically $t$-chromatic, then
\begin{equation}\label{eq:graphsize}
 n\geq \frac12\exp\!\left(\sqrt{(t-1)/\Czero}\right)-1.
\end{equation}
Thus its number of vertices grows superpolynomially in $t$.
\end{corollary}

This contrasts with ordinary chromatic number: there are
triangle-free graphs of chromatic number $t\geq2$ on only
$O(t^2\log t)$ vertices, by Kim's Ramsey bound \cite{Kim}.
We recall this consequence in the graph subsection.

 A partial sign matrix
$A\in\{-1,1,*\}^{M\times N}$ determines a simplicial complex~$S(A)$ with free involution, the sign complex of~$A$: each row specifies a simplex of signed column labels, and
we include these simplices together with their antipodes
\cite{FHV}. Every finite simplicial complex with free involution arises this way,
so the simplicial case of Theorem~\ref{thm:main} is equivalent to
the bound
\[
 \ind|S(A)|\leq C_0\log(2M)\log(2N).
\]
The inequality
$\ind|S(A)|+1\leq\operatorname{signrank}(A)$ makes the index a
topological lower bound for sign-rank~\cite{FHV}. Our theorem
bounds this invariant in terms of the VC dimension as well:
for $N\geq2$ and VC dimension $d\geq1$, we obtain
$O(d\log^2 N)$ for total matrices and $O(d\log^3 N)$ for partial
matrices, independently of the number of rows; see
Corollaries~\ref{cor:total} and~\ref{cor:partial}.


\paragraph{The Morse-theoretic mechanism.}
We encode vertices and facets by a bounded matrix $A$. A smooth
approximation to the $\ell_1$-sphere carries the even function
$x\mapsto\|Ax\|_p^p$. At a high critical value, the curvature of the
constraint contributes a negative definite term to its Hessian.
A trace estimate limits the number of nonnegative directions, and
Morse theory turns this into an index bound for a superlevel set.
Choosing $p\asymp\log f$ and a rounding exponent
$q-1\asymp1/\log v$ yields the two logarithms.

\section{The Morse-theoretic index bound}\label{sec:morse}

\subsection{Preliminary lemmas}

For a symmetric quadratic form $H$, let $\nuplus(H)$ be its number of
nonnegative eigenvalues, including nullity, and let $\nu_{>0}(H)$ count
its positive eigenvalues. These numbers do not depend on the inner
product used to represent the form.
Write $H_1\preceq H_2$ if $H_2-H_1$ is positive semidefinite.

The following lemma combines three standard facts. First, a regular
superlevel set of a Morse function has a CW model whose cell dimensions
are the numbers of positive Hessian eigenvalues at the contributing
critical points
\cite[Theorem~3.5 and the subsequent remark]{Milnor}.
Second, for a free involution, a CW model of the quotient lifts to an
equivariant CW model of the covering space, of the same dimension
\cite[Proposition~1.30 and Appendix, Exercise~1]{Hatcher}.
Third, elementary obstruction theory gives an equivariant map from
any CW complex~$X$ with free involution of dimension at most $D$ to~$S^D$, that is, $\ind X\le \dim X$. Indeed,
choose the map equivariantly on the vertices, then extend over one
cell from each orbit and prescribe the extension on its partner by
equivariance. These extensions exist because $S^D$ is
$(D-1)$-connected; see \cite[Lemma~4.7]{Hatcher}.
A small Morse perturbation allows degenerate critical points as well,
provided the Hessian bound includes nullity.

\begin{lemma}[Superlevel index from Hessian bounds]\label{lem:morse}
Let $\Sigma$ be a compact smooth manifold without boundary, with a
smooth free involution, and let $f\colon\Sigma\to\R$ be smooth and invariant.
Suppose $a<b$ and $D\geq0$ is an integer such that
\[
 \nuplus(\Hess_\Sigma f(x))\leq D
 \quad\text{whenever }d_\Sigma f(x)=0\text{ and }f(x)\geq a.
\]
Then
\[
 \ind\{f\geq b\}\leq D.
\]
\end{lemma}

\begin{proof}
Let $\pi\colon\Sigma\to Q=\Sigma/\mathbb Z_2$ be the quotient covering,
and let $\bar f\colon Q\to\R$ be the descended function.
Since $\pi$ is a local diffeomorphism, the hypothesis says that
$\bar f$ has at least $\dim Q-D$ negative Hessian eigenvalues
at every critical point with value at least $a$.
Put $c=(a+b)/2$.
By Morse approximation and stability of this lower bound under
$C^2$ perturbations
\cite[Corollary~6.8 and Lemma~22.4]{Milnor},
we can choose a Morse function $\bar h$ sufficiently close to $\bar f$
that
\[
 \|\bar h-\bar f\|_\infty<\min\{c-a,b-c\},
\]
and every critical point of $\bar h$ with value at least $c$ has
at most $D$ positive Hessian eigenvalues.
An arbitrarily small constant shift makes $c$ a regular value
while preserving these properties.

Set $B=\{\bar h\geq c\}$ and $W=\pi^{-1}(B)$.
The uniform approximation gives $\{f\geq b\}\subseteq W$.
Applying Morse theory to $-\bar h$, every critical point contributing
to the sublevel set $B=\{-\bar h\leq-c\}$ has Morse index at most $D$.
Thus $B$ has the homotopy type of a finite CW complex $K$ with
$\dim K\leq D$.
Pulling the double cover $W\to B$ back along a homotopy equivalence
$K\to B$ gives a free $\mathbb Z_2$-CW complex~$\widetilde K$,
of the same dimension, equivariantly homotopy equivalent to $W$.
The standard dimension bound for index therefore yields
\[
 \ind\{f\geq b\}
 \leq \ind W
 = \ind\widetilde K
 \leq D.
\]
\end{proof}

Including nullity makes the estimate stable under perturbation.
The gap between $a$ and $b$ permits the perturbation without asserting
that the original, possibly degenerate superlevel is a Morse sublevel.

The following lemma follows by combining Weyl's monotonicity theorem
\cite[Corollary~4.3.3, p.~182]{HornJohnson} with Sylvester's law of inertia
\cite[Theorem~4.5.8, p.~222]{HornJohnson}.
The trace bound then uses the elementary observation that the number of
eigenvalues of a positive semidefinite matrix that are at least~$1$ is
at most its trace. We include a proof for completeness.

\begin{lemma}[Trace bound for nonnegative directions]
\label{lem:trace-bound}
Let $P,D\in\mathbb R^{n\times n}$ be symmetric matrices, with $P$
positive semidefinite and $D$ positive definite, and let $\beta>0$.
Let $H$ be a symmetric bilinear form on a subspace
$E\subseteq\mathbb R^n$, satisfying
\[
 H(v,v)\leq v^{\mathsf T}Pv-\beta v^{\mathsf T}Dv
 \qquad\text{for every }v\in E.
\]
Then
\begin{equation}\label{eq:tracelemma}
 \nu_{\geq0}(H)
 \leq \beta^{-1}\operatorname{tr}(D^{-1/2}PD^{-1/2}).
\end{equation}
\end{lemma}

\begin{proof}
We first identify a subspace on which the positive term $P$ dominates
$\beta D$, then change coordinates so that $D$ becomes the identity.

\smallskip
\noindent\emph{Step 1: isolate the nonnegative directions of $H$.}
Set $k=\nu_{\geq0}(H)$. By the spectral theorem, the eigenvectors
corresponding to the nonnegative eigenvalues of $H$ span a
$k$-dimensional subspace $V\subseteq E$ on which $H$ is nonnegative.
Thus, for every $v\in V$,
\[
 0\leq H(v,v)\leq v^{\mathsf T}Pv-\beta v^{\mathsf T}Dv,
 \qquad\text{so}\qquad
 v^{\mathsf T}Pv\geq\beta v^{\mathsf T}Dv.
\]

\smallskip
\noindent\emph{Step 2: normalize the positive definite form $D$.}
Since $D$ is positive definite, its symmetric positive definite
square root $D^{1/2}$ is invertible. Put
\[
 M=D^{-1/2}PD^{-1/2},\qquad U=D^{1/2}V.
\]
The change of coordinates is invertible, so $\dim U=k$.
For $u=D^{1/2}v\in U$, the inequality from Step~1 becomes
\[
 u^{\mathsf T}Mu
 =v^{\mathsf T}Pv
 \geq\beta v^{\mathsf T}Dv
 =\beta\|u\|_2^2.
\]
Thus every unit vector in $U$ has quadratic value at least $\beta$
under $M$, which is positive semidefinite on all of $\mathbb R^n$.

\smallskip
\noindent\emph{Step 3: sum the contributions to the trace.}
Choose an orthonormal basis $u_1,\ldots,u_k$ of $U$ and extend it
to an orthonormal basis $u_1,\ldots,u_n$ of $\mathbb R^n$.
In this basis, the first $k$ diagonal entries of $M$ are at least
$\beta$ and the remaining entries are nonnegative. Consequently,
\[
 \operatorname{tr}(M)
 =\sum_{i=1}^n u_i^{\mathsf T}Mu_i
 \geq\sum_{i=1}^k u_i^{\mathsf T}Mu_i
 \geq k\beta.
\]
Dividing by $\beta$ gives the claimed bound.
\end{proof}

\subsection{The matrix bound}

The matrix formulation below records the vertex--facet incidences of a
free simplicial complex. Choose one vertex from each antipodal pair and
realize these vertices and their partners as $e_j$ and $-e_j$ on the
$\ell_1$-sphere. A row of the signed incidence matrix $A$ records which
of these vertices belong to a chosen facet: its entries are $1$, $-1$,
or $0$. On that facet the corresponding coordinate of $Ax$ is $1$,
and on its antipodal facet it is $-1$. Since $\|Ax\|_\infty\leq
\|x\|_1=1$, the realization lies in the equality set $R(A)$ defined
below. Thus the numbers of vertex and facet orbits become the numbers
of columns and rows. Section~\ref{sec:complexes} gives the full reduction,
including the extension to strongly regular CW complexes.

Gromov's nonlinear spectra provide a methodological precedent: they
measure energy sublevel sets by topological invariants rather than
counting linear eigenvalues \cite[\S\S0.3--0.4]{Gromov}.
For the norm ratio $\|x\|_p/\|x\|_q$ on a finite set with uniform
measure, he identifies the critical directions, whose nonzero
coordinates have equal absolute values, and computes their Morse
indices \cite[\S1.3.A, p.~150]{Gromov}. Up to smoothing, our energy is
$E_A(x)=\|Ax\|_p/\|x\|_q$; the case $A=I$ recovers this model,
up to the constant arising from the normalization of the measure.
Operator norm ratios also appear in his broader framework
\cite[\S3.1, pp.~153--154]{Gromov}. For general bounded matrices,
we replace the classification of critical points by a uniform trace
bound on the number of nonnegative Hessian eigenvalues, including
nullity. Gromov's essential dimension measures deformation into
low-dimensional subsets of an ambient projective space; our argument
instead constructs a low-dimensional CW model of a containing
superlevel set after perturbation. Lifting this model to its double
cover bounds the ordinary $\mathbb Z_2$-index.

\begin{theorem}[Matrix bound]\label{thm:matrix}
For $m,n\geq1$ and $A\in[-1,1]^{m\times n}$, set
\[
 R(A)=\{x\in\R^n:\|x\|_1=\|Ax\|_\infty=1\}.
\]
With the antipodal involution,
\begin{equation}\label{eq:matrix}
 \ind R(A)\leq
 \min\{m-1,n-1,\Czero\log(2m)\log(2n)\}.
\end{equation}
\end{theorem}

The proof replaces the maximum over the $m$ rows by an $\ell_p$-norm
and the $\ell_1$-sphere in $\R^n$ by a smoothed $\ell_q$-sphere.
Large $p$ approximates the
maximum, whereas $q$ close to $1$ approximates the original constraint.
Both choices have a curvature cost. The Hessian estimate contains the
ratio $(p-1)/(q-1)$, comparing the positive contribution of the objective
with the negative contribution of the constraint, and the squared
norm-comparison losses $m^{2/p}n^{2-2/q}$; see
\eqref{eq:central}. Increasing $p$ or decreasing $q$ improves the norm
comparison but increases the curvature ratio. Taking
$p\asymp\log(2m)$ and $q-1\asymp1/\log(2n)$ keeps both dimension
factors bounded, leaving a curvature ratio of order~$\log(2m)\log(2n)$.

\begin{proof}
The
odd maps $x\mapsto Ax/\|Ax\|_2$ and $x\mapsto x/\|x\|_2$ give
the bounds $m-1$ and $n-1$, respectively, and settle $\min(m,n)=1$.
Assume henceforth that $m,n\geq2$.
We map $R(A)$ into a smooth superlevel set, bound the Hessian
at every sufficiently high critical point, and apply
Lemma~\ref{lem:morse}.

\smallskip
\noindent\emph{Step 1: choose a smooth constraint.}
Let $p\geq4$ be an even integer and let $1<q<2$.
We leave these exponents unspecified until the final estimate.
Define
\[
 G(x)=\sum_{j=1}^n\bigl(x_j^2+(4n)^{-2}\bigr)^{q/2},\qquad
 \Sigma=\{G=1\},\qquad f(x)=\|Ax\|_p^p.
\]
The fixed positive term $(4n)^{-2}$ smooths the zero coordinates while
keeping their total contribution $G(0)=n(4n)^{-q}$ below $1/4$.
This leaves a fixed margin for the radial comparison in Step~2;
no limiting argument in the smoothing parameter is needed.
The function $G$ is proper, and
\[
 \frac{\partial G}{\partial x_j}(x)
 =qx_j\bigl(x_j^2+(4n)^{-2}\bigr)^{q/2-1}.
\]
Thus $\nabla G$ vanishes only at $0$, which is not on $\Sigma$.
The level set $\Sigma$ is therefore a compact smooth hypersurface
with a free antipodal involution. Taking $p$ even makes
$f(x)=\sum_i(Ax)_i^p$ a smooth even polynomial, including where
some $(Ax)_i=0$. The convenient restriction $p\geq4$ also avoids
an endpoint case in the first use of H\"older's inequality in Step~5.

\smallskip
\noindent\emph{Step 2: send $R(A)$ into a superlevel.}
For $u\ne0$, the function $s\mapsto G(su)$ is strictly increasing
for $s>0$, starts below $1/4$, and tends to infinity.
There is a unique $s(u)>0$ with $G(s(u)u)=1$.
The implicit function theorem gives continuity of~$s$, and evenness of~$G$
gives $s(-u)=s(u)$.

If $\|u\|_1=1$, subadditivity of $t^{q/2}$ and
$|u_j|^q\leq|u_j|$ give
\[
 1=G(s(u)u)
 \leq s(u)^q\sum_j|u_j|^q+n(4n)^{-q}
 \leq s(u)^q+\tfrac14,
\]
so $s(u)\geq(3/4)^{1/q}>3/4$.
For $u\in R(A)$, we have $\|Au\|_\infty=1$, and hence
\begin{equation}\label{eq:levels}
 f(s(u)u)=s(u)^p\|Au\|_p^p
 \geq s(u)^p>(3/4)^p>2^{-p}.
\end{equation}
The odd radial map $u\mapsto s(u)u$ therefore sends $R(A)$ into
$\{x\in\Sigma:f(x)\geq(3/4)^p\}$.
By Lemma~\ref{lem:morse}, it suffices to bound the nonnegative
Hessian count at critical points $z$ with $f(z)\geq2^{-p}$.
The levels $2^{-p}$ and $(3/4)^p$ serve two purposes: their gap
leaves room for the Morse perturbation, and the lower level makes
the factor $f(z)^{-2/p}$ in the final estimate at most $4$.


\smallskip
\noindent\emph{Step 3: bounding the Hessian of~$f$.}
Restricting $f$ to the curved hypersurface $\Sigma$ introduces a
negative term into its Hessian. We compute this term and bound its
size in five parts.

Fix a critical point $z\in\Sigma$ of $f|_\Sigma$ with
$f(z)\geq2^{-p}$. This point remains fixed throughout Steps~3--5.
Set
\[
 r_j=\bigl(z_j^2+(4n)^{-2}\bigr)^{1/2},\qquad
 D=\diag(r_j^{q-2}),\qquad
 W=\diag(|(Az)_i|^{p-2}).
\]
Here $D$ is positive definite and $W$ is positive semidefinite.
All three quantities are evaluated at the fixed critical point~$z$.

\smallskip
\noindent\emph{(a) Compute the ambient Hessian of $f$.}
For any $v\in\R^n$, evenness of $p$ gives
\[
 f(z+tv)=\sum_i\bigl((Az)_i+t(Av)_i\bigr)^p.
\]
Differentiating twice at $t=0$, we obtain
\[
 \nabla^2f(z)[v,v]
 =p(p-1)\sum_i|(Az)_i|^{p-2}(Av)_i^2.
\]
Thus the ambient Hessian is positive semidefinite and, in matrix form,
\[
 \nabla^2f(z)=p(p-1)A^TWA.
\]

\smallskip
\noindent\emph{(b) Bound the curvature of the constraint.}
The mixed second derivatives of $G$ vanish, and differentiation gives
\begin{align*}
 \frac{\partial G}{\partial x_j}(z)
 &=qz_jr_j^{q-2},\\
 \frac{\partial^2G}{\partial x_j^2}(z)
 &=qr_j^{q-2}+q(q-2)z_j^2r_j^{q-4}.
\end{align*}
The second term is nonpositive because $q<2$.
Substituting $z_j^2=r_j^2-(4n)^{-2}$ reveals a positive lower bound:
\begin{align*}
 \frac{\partial^2G}{\partial x_j^2}(z)
 &=q(q-1)r_j^{q-2}
   +\frac{q(2-q)}{(4n)^2}r_j^{q-4}\\
 &\geq q(q-1)r_j^{q-2}.
\end{align*}
Consequently, for every $v\in\R^n$,
\[
 \nabla^2G(z)[v,v]
 \geq q(q-1)\sum_jr_j^{q-2}v_j^2.
\]
Equivalently,
\[
 \nabla^2G(z)\succeq q(q-1)D.
\]

\smallskip
\noindent\emph{(c) Estimate the Lagrange multiplier.}
The tangent space is
\[
 T_z\Sigma=\{v:\langle\nabla G(z),v\rangle=0\}.
\]
Criticality of $f|_\Sigma$ says that $\nabla f(z)$ is also
orthogonal to this tangent space. Since $\nabla G(z)\ne0$,
there is a scalar $\lambda$ such that
\[
 \nabla f(z)=\lambda\nabla G(z).
\]
To determine its sign and size, take the inner product with $z$.
Differentiating $f(tz)=t^pf(z)$ at $t=1$ gives
\[
 \langle z,\nabla f(z)\rangle=pf(z).
\]
On the other hand,
\[
 0<\langle z,\nabla G(z)\rangle
 =q\sum_jz_j^2r_j^{q-2}
 \leq q\sum_jr_j^q=q.
\]
Here we used $z\ne0$, $z_j^2\leq r_j^2$, and $G(z)=1$.
Therefore
\[
 \lambda=\frac{pf(z)}{\langle z,\nabla G(z)\rangle}
 \geq\frac{pf(z)}q>0.
\]

\smallskip
\noindent\emph{(d) Compute the Hessian of $f$ on $\Sigma$.}
For $v\in T_z\Sigma$, choose a smooth curve $\gamma$ in $\Sigma$
with $\gamma(0)=z$ and $\gamma'(0)=v$.
The chain rule gives
\[
 (f\circ\gamma)''(0)
 =\nabla^2f(z)[v,v]
  +\langle\nabla f(z),\gamma''(0)\rangle.
\]
Since $G(\gamma(t))=1$, differentiating twice gives
\[
 0=\nabla^2G(z)[v,v]
   +\langle\nabla G(z),\gamma''(0)\rangle.
\]
Substituting $\nabla f(z)=\lambda\nabla G(z)$ into the first
identity gives
\begin{align*}
 (f\circ\gamma)''(0)
 &=\nabla^2f(z)[v,v]
   +\lambda\langle\nabla G(z),\gamma''(0)\rangle\\
 &=\nabla^2f(z)[v,v]-\lambda\nabla^2G(z)[v,v].
\end{align*}
At a critical point, this second derivative is the constrained
Hessian $H(v,v)$, where $H=\Hess_\Sigma f(z)$; it depends only
on the tangent vector $v$, not on the chosen curve. Hence
\[
 H=\bigl(\nabla^2f(z)-\lambda\nabla^2G(z)\bigr)|_{T_z\Sigma}.
\]

\smallskip
\noindent\emph{(e) Combine the estimates.}
For every tangent vector $v$, parts (b) and (c) give
\[
 \lambda\nabla^2G(z)[v,v]
 \geq\frac{pf(z)}q\,q(q-1)\sum_jr_j^{q-2}v_j^2
 =p(q-1)f(z)\sum_jr_j^{q-2}v_j^2.
\]
Subtract this lower bound from the ambient Hessian in part (a),
and divide by $p$. We obtain
\[
 \frac1pH(v,v)
 \leq(p-1)\sum_i|(Az)_i|^{p-2}(Av)_i^2
      -(q-1)f(z)\sum_jr_j^{q-2}v_j^2.
\]
In matrix notation, this is
\begin{equation}\label{eq:hessian}
 \frac1p H\preceq
 \bigl((p-1)A^TWA-(q-1)f(z)D\bigr)|_{T_z\Sigma}.
\end{equation}

\smallskip
\noindent\emph{Step 4: bound the number of nonnegative directions by trace.}
Apply Lemma~\ref{lem:trace-bound} to \eqref{eq:hessian}.
Multiplication by $1/p>0$ does not change the Hessian count, so
\begin{align}
 \nuplus(H)
 &\leq\frac{p-1}{(q-1)f(z)}
       \tr(D^{-1/2}A^TWAD^{-1/2})\notag\\
 &=\frac{p-1}{(q-1)f(z)}
       \sum_{i,j}|(Az)_i|^{p-2}A_{ij}^2r_j^{2-q}\notag\\
 &\leq\frac{p-1}{(q-1)f(z)}
       \left(\sum_i|(Az)_i|^{p-2}\right)
       \left(\sum_jr_j^{2-q}\right).
 \label{eq:traceproduct}
\end{align}
The equality computes diagonal entries, using
$(D^{-1})_{jj}=r_j^{2-q}$; the last inequality uses $A_{ij}^2\leq1$.

\smallskip
\noindent\emph{Step 5: bound the two sums and apply Lemma~\ref{lem:morse}.}
H\"older's inequality, applied to each sequence and the constant
sequence $1$, gives
\begin{align*}
 \sum_i|(Az)_i|^{p-2}
 &\leq m^{2/p}\left(\sum_i|(Az)_i|^p\right)^{1-2/p}
 =m^{2/p}f(z)^{1-2/p},\\
 \sum_jr_j^{2-q}
 &\leq n^{2-2/q}\left(\sum_jr_j^q\right)^{(2-q)/q}
 =n^{2-2/q}.
\end{align*}
The last equality uses $G(z)=1$.
Substituting into \eqref{eq:traceproduct} gives
\begin{equation}\label{eq:central}
 \nuplus(H)\leq
 \frac{p-1}{q-1}\,m^{2/p}n^{2-2/q}f(z)^{-2/p}.
\end{equation}
We now choose
\[
 p=2\lceil\log(2m)\rceil,\qquad
 q=1+\frac1{2\log(2n)}.
\]
These satisfy the required conditions $p\geq4$ even and $1<q<2$.
Their logarithmic scales keep the two dimension factors bounded:
\[
 m^{2/p}
 \leq\exp\!\left(\frac{\log m}{\log(2m)}\right)\leq e,
 \qquad
 n^{2-2/q}
 =\exp\!\left(\frac{\log n}{\log(2n)+1/2}\right)\leq e.
\]
Also $p-1\leq2\log(2m)+1\leq3\log(2m)$, since
$\log(2m)\geq\log4>1$. Thus
\[
 \frac{p-1}{q-1}
 =2(p-1)\log(2n)\leq6\log(2m)\log(2n).
\]
This quotient produces the two logarithms.
Finally, $f(z)\geq2^{-p}$ gives $f(z)^{-2/p}\leq4$, so
\eqref{eq:central} yields
\[
 \nuplus(H)\leq24e^2\log(2m)\log(2n)
 =\Czero\log(2m)\log(2n).
\]
The critical point $z$ was arbitrary subject to $f(z)\geq2^{-p}$,
so this bound holds at every critical point required by
Lemma~\ref{lem:morse}.
Since the Hessian count is an integer, that lemma applies with
the integer part of this bound and the levels $2^{-p}<(3/4)^p$.
Composing with the odd map from Step~2 gives
\[
 \ind R(A)
 \leq\ind\{x\in\Sigma:f(x)\geq(3/4)^p\}
 \leq\Czero\log(2m)\log(2n),
\]
as claimed.
\end{proof}

\begin{remark}[Robust matrix bound]\label{rem:robust-matrix}
More generally, for $0<\rho\leq1$ and $A\in[-1,1]^{m\times n}$, let
\[
 R_\rho(A)=\{x\in\R^n:\|x\|_1=1,\ \|Ax\|_\infty\geq\rho\}.
\]
Then
\[
 \ind R_\rho(A)\leq
 \min\{m-1,n-1,\Czero\rho^{-2}\log(2m)\log(2n)\}.
\]
The applications below use only $\rho=1$.
\end{remark}

\begin{remark}[VC Refinement for Total Matrices]\label{rem:refinement}
For a total sign matrix $A\in\{-1,+1\}^{m\times n}$ with row
VC dimension $d\geq1$, we can refine the bound to
\[
 \operatorname{ind} R(A)
 \leq C_d\bigl(\log(2m)\log(2n)\bigr)^{d/(d+1)},
\]
where $C_d$ depends only on $d$.
The idea is to replace the Hessian trace estimate by an eigenvalue
counting estimate. It is possible to put the normalized positive matrix in the Hessian
calculation into the form
\[
 M=\sum_i\pi_i b_i b_i^{\mathsf T},
 \qquad b_i=(A_{ij}\sqrt{\mu_j})_{j=1}^n,
\]
where $\pi$ and $\mu$ are some probability distributions on the rows and
columns, respectively. Haussler's packing theorem~\cite{Haussler1995}
provides $O_d(\eta^{-d})$ representative rows such that every row
disagrees with a representative on columns of total $\mu$-mass
at most $\eta$. The corresponding weighted vectors therefore
satisfy $\lVert b_i-b_{i'}\rVert_2^2\leq4\eta$.

Let $V$ be the linear span of these representative weighted vectors.
Then $\dim V=O_d(\eta^{-d})$, and every $b_i$ has squared distance
at most $4\eta$ from $V$. Thus, writing $P_V$ for orthogonal
projection onto $V$, we have
$\operatorname{tr}((I-P_V)M)\leq4\eta$.
It follows that $M$ has at most $\dim V+4\eta/t$ eigenvalues
at least $t$. Choosing $\eta=t^{1/(d+1)}$ gives
$O_d(t^{-d/(d+1)})$. Substituting this estimate into the
constrained Hessian calculation and applying the same
Morse-theoretic argument proves the refinement.
\end{remark}


\section{Free complexes and the FLM inequality}\label{sec:complexes}

\begin{proof}[Proof of Theorem~\ref{thm:main}]
First let $K$ be a finite simplicial complex with free involution with $v=2n$ vertices
and $f=2m$ facets. No simplex contains an antipodal vertex pair, since
the midpoint of such a pair would be fixed. Choose representatives
$v_1,\ldots,v_n$ of the vertex orbits and
$\sigma_1,\ldots,\sigma_m$ of the facet orbits. Set
\[
 A_{ij}=\begin{cases}
 1,&v_j\in\sigma_i,\\
 -1,&\tau v_j\in\sigma_i,\\
 0,&\text{otherwise}.
 \end{cases}
\]
Realize $v_j$ as $e_j$ and $\tau v_j$ as $-e_j$ on the
$\ell_1$-sphere. On $\sigma_i$, the $i$th coordinate of $Ax$ equals
$1$, and on $\tau\sigma_i$ it equals $-1$. Hence
$|K|\subseteq R(A)$, and Theorem~\ref{thm:matrix} gives
\[
 \ind|K|\leq\Czero\log(2n)\log(2m)=\Czero\log v\log f.
\]

Now let $X$ be strongly regular. Every closed cell $\sigma$ misses
$\tau\sigma$: a nonempty intersection would be an invariant closed
cell, hence a ball on which the involution has a fixed point by
Brouwer's theorem. In particular, no closed cell contains antipodal
vertices.

Form a simplicial complex $K$ by placing a full simplex on the vertex
set of each facet of $X$. It has $v$ vertices, at most $f$ facets, and
a free involution. There is an equivariant map $X\to|K|$ carrying each
closed cell into the simplex on its vertices. To construct it, start
with the vertex map and extend over one cell in each orbit. Its
boundary has already mapped into that simplex, which is contractible.
Define the extension on the paired cell by equivariance.
Monotonicity of the index and the simplicial case complete the proof.
\end{proof}

\begin{remark}[Scope and sharpness]
\emph{Comparison with Cohen--Macaulay complexes.}
For a centrally symmetric simplicial $(d-1)$-sphere, the stronger
facet bound $f\geq2^d$ is classical. Stanley~\cite[Theorem~3.1]{Stanley} proves it more generally
for centrally symmetric simplicial complexes that are Cohen--Macaulay
over~$\R$. Theorem~\ref{thm:main} applies without this
homological hypothesis and includes nonsimplicial cell structures.
The lifted sphere in Corollary~\ref{cor:rp} may have exponentially many
facets, so this classical facet bound does not yield a superpolynomial
vertex bound. The enlargement used in its proof has few facets but need not be Cohen--Macaulay.

\smallskip
\noindent\emph{Sharpness of the logarithmic product.}
The product-of-logarithms order is sharp:
centrally symmetric Hanner polytopes can have
$\log v,\log f=\Theta(\sqrt d)$; see~\cite{FLM} and~\cite[Proposition~4.1]{Milo}.
The preceding simplicial enlargement retains an equivariant map from
$S^{d-1}$ and does not increase either count, so it gives the same
sharpness for arbitrary free simplicial complexes.

\smallskip
\noindent\emph{Necessity of strong regularity.}
The antipodal CW structure on $S^d$
with two cells in each dimension is regular and has two vertices and
two facets, but its index is $d$; for $d\geq1$ its closed cells do not
satisfy the intersection condition.
\end{remark}

\section{The size consequences}

\subsection{Projective-space triangulations}\label{subsection:triangulations}

\begin{proof}[Proof of Corollary~\ref{cor:rp}]
Lift a triangulation $T$ of $\RP^d$ with $N$ vertices along
$S^d\to\RP^d$. The resulting simplicial complex $K$ has $2N$ vertices,
and its deck involution makes $|K|$ equivariantly homeomorphic to $S^d$.

For $w\in V(K)$, let $N[w]$ be its closed neighborhood in the
$1$-skeleton. No $N[w]$ contains both $v$ and $\tau v$.
There is no edge from $v$ to $\tau v$, since an edge maps injectively
to a base edge. If $w\notin\{v,\tau v\}$ and both $v,\tau v\in N[w]$,
the two edges $wv$ and $w\tau v$ would lift the same base edge with
the same initial vertex, contradicting uniqueness of path lifting.

Let $2^U$ denote the full simplex on a vertex set $U$, and set
\[
 L=\bigcup_{w\in V(K)}2^{N[w]}.
\]
The involution preserves $L$ and is free, since no generating simplex
contains an antipodal pair. Moreover, $K\subseteq L$: the vertices
of a simplex lie in the closed neighborhood of any of its vertices.
The complex $L$ has $2N$ vertices and at most $2N$ facets. Hence
\[
 d=\ind|K|\leq\ind|L|\leq\Czero\log^2(2N). \qedhere
\]
\end{proof}

Cup products already give vertex lower bounds without a manifold
hypothesis. In particular, the estimates of Govc, Marzantowicz,
and Pave\v{s}i\'c \cite[Theorems~1.2 and~3.5]{GMP} imply that
$\alpha^h\ne0$, for a degree-one cohomology class on an
$N$-vertex simplicial complex, forces
$N\geq(h+1)(h+2)/2$. For powers of a single degree-one
mod-two class, the following corollary gives instead the
superpolynomial bound $N\geq\frac12\exp(\sqrt{h/\Czero})$.

\begin{corollary}[Degree-one cup powers]\label{cor:cup}
Let $T$ be a finite simplicial complex with $N$ vertices.
If $\alpha\in H^1(T;\mathbb F_2)$ satisfies $\alpha^h\ne0$ for an
integer $h\geq1$, then
\begin{equation}\label{eq:cup}
 h\leq\Czero\log^2(2N).
\end{equation}
In particular, \eqref{eq:rp} holds for every finite simplicial complex
homotopy equivalent to $\RP^d$, $d\geq1$.
\end{corollary}

\begin{proof}
Let $K\to T$ be the double cover classified by $\alpha$.
An equivariant map $|K|\to S^k$ induces $|T|\to\RP^k$ pulling
the class of the double cover $S^k\to\RP^k$ back to $\alpha$.
Its $(k+1)$st power vanishes, since $H^{k+1}(\RP^k;\mathbb F_2)=0$.
Thus $\alpha^h\ne0$ forces $k\geq h$, and $\ind|K|\geq h$.
The neighborhood enlargement from the proof of Corollary~\ref{cor:rp}
supplies the upper bound. For the last assertion, use the degree-one
class whose $d$th power is nonzero in $H^*(\RP^d;\mathbb F_2)$.
\end{proof}

In the language of covering type, Corollary~\ref{cor:cup} gives
\[
 \operatorname{ct}(\RP^d)
 \geq \frac12\exp\!\left(\sqrt{d/\Czero}\right).
\]
Here $\operatorname{ct}(X)$, introduced by Karoubi and Weibel
\cite{KW}, is the minimum size of a good open cover of a space
homotopy equivalent to $X$; a cover is good if every nonempty
finite intersection of its members is contractible.
For spaces of finite simplicial homotopy type, this equals the
minimum number of vertices in a finite simplicial complex
homotopy equivalent to $X$ \cite[Theorem~1.2]{GMP}.

\subsection{Triangle-free graphs}\label{sec:graphs}

For a finite simple graph $G$ on $V$, write $v^+=(v,+)$ and
$v^-=(v,-)$. The complex $B_0(G)$ consists of all simplices
\[
 A^+\cup B^-,\qquad A,B\subseteq V,\quad A\cap B=\varnothing,
 \quad ab\in E(G)\text{ for every }a\in A,\ b\in B.
\]
Either shore may be empty. Exchanging signs is a free involution.
The standard box-complex bound is
\[
 \chi(G)\geq\ind|B_0(G)|+1
 \geq\coind|B_0(G)|+1,
\]
as recalled in \cite{ST}. Thus the definition in the introduction
is equivalent to $\coind|B_0(G)|\geq t-1$.

\begin{proof}[Proof of Corollary~\ref{cor:graph}]
For each $v\in V$, let $N(v)$ be its open neighborhood and put
\[
 F_v=N(v)^+\cup(V\setminus N(v))^-.
\]
Let $L$ be generated by the simplices $F_v$, their antipodes, and
the two pure shores $V^+$ and $V^-$. Each generating simplex contains
exactly one of $w^+,w^-$ for every $w\in V$, so the involution is free.
The complex has $2n$ vertices and at most $2n+2$ facets.

We claim that $B_0(G)\subseteq L$. Pure shores are already included.
For a mixed simplex $A^+\cup B^-$, choose $b\in B$.
Then $A\subseteq N(b)$. Any edge within $B$, together with any
$a\in A$, would form a triangle. Hence $B$ is independent and
$B\subseteq V\setminus N(b)$.
Thus $A^+\cup B^-\subseteq F_b$, proving the claim.

Theorem~\ref{thm:main} now yields \eqref{eq:graphindex}.
If $G$ is topologically $t$-chromatic, monotonicity gives
\[
 t-1\leq\ind|B_0(G)|
 \leq\Czero\log(2n)\log(2n+2)
 \leq\Czero\log^2(2n+2).
\]
Rearranging proves \eqref{eq:graphsize}.
\end{proof}

\begin{remark}[Chromatic number and topological bounds]
For every integer $c\geq2$, there are triangle-free graphs with ordinary
chromatic number $c$ on only $O(c^2\log c)$ vertices.
Indeed, Kim's Ramsey bound \cite{Kim} gives $n$-vertex triangle-free
graphs with independence number $\alpha(G)=O(\sqrt{n\log n})$.
Using $\chi(G)\geq n/\alpha(G)$, choose $n=O(c^2\log c)$ so that
$\chi(G)\geq c$. Deleting one vertex decreases the chromatic number
by at most one, so successive deletions give $\chi(G)=c$.
For these graphs, Corollary~\ref{cor:graph} gives
$\ind|B_0(G)|=O(\log^2 c)$.
\end{remark}

\subsection{Sign complexes}

Let $A \in \{-1,1,*\}^{M \times N}$ be a partial sign matrix with
columns indexed by $[N]$. Following \cite[Section~3.1]{FHV}, its
\emph{sign complex} $S(A)$ is the simplicial complex generated by
the following simplices on the signed column set
$\{1^{+},1^{-},\dots,N^{+},N^{-}\}$:
\[
  \sigma_r^{+}
  = \{\, i^{+} : r_i = +1 \,\}
    \cup \{\, i^{-} : r_i = -1 \,\},
  \qquad
  \sigma_r^{-} = -\,\sigma_r^{+},
\]
as $r$ ranges over the rows of $A$. Exchanging $i^{+}$ and $i^{-}$
defines a free involution. The complex has at most $2N$ vertices
and at most $2M$ facets; a column consisting entirely of $*$ entries
contributes no vertices. If the realization is nonempty,
Theorem~\ref{thm:main} gives the following bound. 

\begin{corollary}\label{cor:sign}
For every partial sign matrix $A \in \{-1,1,*\}^{M\times N}$,
\[
  \operatorname{ind} |S(A)|
  \;\le\; C_0 \log(2N)\log(2M).
\]
In particular $\operatorname{ind}|S(A)| \le C_0 \log^2(2N)$
for a square matrix.
\end{corollary}
Every finite simplicial complex with a free involution arises as the sign complex of a partial matrix \cite[Lemma~15]{FHV}, so the simplicial case of
Theorem~\ref{thm:main} and the preceding statement are two readings
of the same result. Given a partial sign matrix $A$, let $\operatorname{VC}(A)$ denote the VC dimension of $A$: the largest $d$ such that some $d$ columns of $A$ carry every pattern in $\{-1, 1\}^d$ among the rows of $A$ restricted to these $d$ columns. If $A \in \{-1, 1\}^{M \times N}$, we call $A$ a total sign matrix.

\begin{corollary}\label{cor:total}
Let $A$ be a total sign matrix with $N \ge 2$ columns and
$\operatorname{VC}(A) = d \ge 1$.
Then
\[
  \operatorname{ind} |S(A)|
  \;\le\; C_0 \, d \log(2eN/d) \log(2N)
  \;=\; O(d \log^2 N),
\]
irrespective of the number of rows.
\end{corollary}

\begin{proof}
Rows with equal sign patterns give equal simplices, so we may
assume the rows are distinct. By the Sauer--Shelah--Perles lemma
\cite{Sauer},
$M \le \sum_{i \le d}\binom{N}{i} \le (eN/d)^d$ and so
$2M \le (2eN/d)^{d}$ since $d \ge 1$.
Now apply Corollary~\ref{cor:sign}.
\end{proof}

For comparison, Alon, Moran, and Yehudayoff \cite{AMY} construct
total $N\times N$ matrices of VC dimension $2$ whose sign-rank
is of order $\sqrt N$ up to logarithmic factors, whereas
Corollary~\ref{cor:total} bounds their index by $O(\log^2 N)$.
Corollary~\ref{cor:total} gives a bound toward Question~9 of \cite{FHV},
which asks whether $\operatorname{coind}|S(A)|$ is bounded by a
function of $\operatorname{VC}(A)$ alone for total matrices.

Chornomaz, Moran, and Waknine~\cite{chornomaz2025spherical} previously posed an equivalent
question for spherical dimension.
Let $H$ be the row class of $A$, put $H^{\pm}=H\cup(-H)$,
and let $\Delta(H)$ be the union of its row simplices before
adjoining antipodes. Spherical dimension~$\operatorname{sd}(H)$ is the coindex of
$\Delta(H)\cap(-\Delta(H))$, whereas $|S(A)|$ is
$\Delta(H)\cup(-\Delta(H))$. Hence
\[
 \operatorname{coind}|S(A)|=\operatorname{sd}(H^{\pm}),
 \qquad
 \operatorname{VC}(H^{\pm})\leq2\operatorname{VC}(H)+1,
\]
where the VC bound follows from the Sauer--Shelah--Perles lemma.
Conversely, $\Delta(H)\cap(-\Delta(H))\subseteq|S(A)|$.
Thus a bound by a function of VC dimension for either invariant
would give one for the other.

Since $\operatorname{coind}\leq\operatorname{ind}$,
Corollary~\ref{cor:total} gives
$\operatorname{coind}|S(A)|=O(d\log^2 N)$.
Removing the dependence on $N$ remains open.


Blondal, Hatami, Hatami, Lalov, and Tretiak ask the stronger
question of whether the index itself is bounded by a function of
VC dimension \cite[Problem~1.18]{BHHHT}.
They also conjecture a product-of-logarithms bound for the
list-replicability number $\operatorname{LR}(A)$
\cite[Conjecture~1.17]{BHHHT}. Their inequality
$\operatorname{ind}|S(A)|\leq2\operatorname{LR}(A)-1$
\cite[Theorem~1.1]{BHHHT} shows that this conjecture would imply
the total-matrix case of Corollary~\ref{cor:sign}.

\paragraph{Partial matrices.}
For partial matrices, VC dimension does not control the number of
distinct rows: the $2^N$ rows in $\{+1,*\}^{N}$ have VC dimension
$0$.
A total sign matrix $\bar A$ \emph{disambiguates} $A$ if every row
$r$ of $A$ is extended by a row $\bar r$ of $\bar A$, meaning that
$\bar r_i=r_i$ whenever $r_i\ne *$.
The quasipolynomial Sauer--Shelah--Perles lemma controls the number
of such extensions and costs one additional logarithm in the
index bound.

\begin{corollary}\label{cor:partial}
There is an absolute constant $C>0$ such that every partial sign
matrix $A$ with $N\geq2$ columns and $\operatorname{VC}(A)=d$ satisfies
\[
 \operatorname{ind}|S(A)|\leq C d\log^3 N.
\]
\end{corollary}

\begin{proof}
If $d=0$, all specified entries in any fixed column agree.
Consequently, a single total row extends every row of $A$.
The complex $S(A)$ therefore lies in two disjoint opposite simplices,
and~$\ind|S(A)|\leq0$.

Extending a row enlarges its simplex, so
$S(A)\subseteq S(\bar A)$ for every disambiguation $\bar A$.
The quasipolynomial Sauer--Shelah--Perles lemma
\cite[Theorem~12]{AHHM} gives a disambiguation with at most
$N^{cd\log N}$ rows, for an absolute constant $c>0$.
By Corollary~\ref{cor:sign} and monotonicity,
\[
 \operatorname{ind}|S(A)|
 \leq\operatorname{ind}|S(\bar A)|
 \leq C_0\bigl(\log 2+cd\log^2 N\bigr)\log(2N)
 \leq C d\log^3 N,
\]
after increasing the absolute constant $C$.
\end{proof}

\begin{remark}
The quasipolynomial bound on the number of rows in a disambiguation
is nearly optimal. There are partial sign matrices of VC dimension
$1$ for which every disambiguation has at least
$N^{(\log N)^{1-o(1)}}$ rows \cite[Theorem~11]{AHHM}.
This concerns the size of disambiguations and does not establish
sharpness of the index bound in Corollary~\ref{cor:partial}.
\end{remark}

\begin{remark}
Already for partial sign matrices of VC dimension $1$, there are
examples~\cite[Proposition~8]{FHV} with
\[
 \operatorname{ind}|S(A)|\geq\operatorname{coind}|S(A)|
 =\Omega\!\left(\frac{\log^2 N}{\log\log N}\right).
\]
Thus the largest possible index for
VC dimension $1$ is bounded below by
$\Omega(\log^2 N/\log\log N)$ and above by $O(\log^3 N)$.
In particular, dependence on $N$ cannot be eliminated for partial
matrices.
\end{remark}

\paragraph{AI use statement.}
The authors developed the topological Figiel--Lindenstrauss--Milman
approach to lower bounds for the vertex numbers of triangulations
of real projective spaces. LLMs assisted with proof discovery (in particular the use of Lemmas~\ref{lem:morse} and~\ref{lem:trace-bound}), proof details, and
manuscript drafts. The authors subsequently substantially
simplified the proofs and revised the manuscript.

\paragraph{Acknowledgements.} 
FF thanks Dan Guyer, Amzi Jeffs, Eran Nevo, and many other colleagues
for helpful conversations about triangulations of real projective spaces.
FF was supported NSF CAREER Grant DMS 2042428.
EM was supported by the National Science Foundation Graduate Research Fellowship Program under Grant Numbers DGE2140739 and DGE2631988. Any opinions, findings, and conclusions or recommendations expressed in this material are those of the authors and do not necessarily reflect the views of the National Science Foundation.

\begin{thebibliography}{12}

\bibitem{AAK}
K.~Adiprasito, S.~Avvakumov, and R.~Karasev,
\emph{A subexponential size triangulation of $\mathbb RP^n$},
Combinatorica \textbf{42} (2022), 1--8.
\href{https://doi.org/10.1007/s00493-021-4602-x}{doi:10.1007/s00493-021-4602-x}.

\bibitem{AHHM}
N.~Alon, S.~Hanneke, R.~Holzman, and S.~Moran,
\emph{A theory of PAC learnability of partial concept classes},
62nd IEEE Annual Symposium on Foundations of Computer Science
(FOCS 2021), IEEE, 2022, 658--671.
\href{https://doi.org/10.1109/FOCS52979.2021.00070}{doi:10.1109/FOCS52979.2021.00070}.
Full version:
\href{https://arxiv.org/abs/2107.08444}{arXiv:2107.08444}.

\bibitem{AMY}
N.~Alon, S.~Moran, and A.~Yehudayoff,
\emph{Sign rank versus VC dimension},
Proceedings of the 29th Annual Conference on Learning Theory,
Proc. Mach. Learn. Res. \textbf{49} (2016), 47--80.
\href{https://proceedings.mlr.press/v49/alon16.html}{PMLR 49}.

\bibitem{AM}
P.~Arnoux and A.~Marin,
\emph{The K\"uhnel triangulation of the complex projective plane from the
view point of complex crystallography. II},
Mem. Fac. Sci. Kyushu Univ. Ser. A \textbf{45} (1991), 167--244.

\bibitem{BHHHT}
A.~Blondal, H.~Hatami, P.~Hatami, C.~Lalov, and S.~Tretiak,
\emph{Sign-rank, index, and list replicability: Connections and separations},
preprint (2026),
\href{https://arxiv.org/abs/2606.18236}{arXiv:2606.18236}.

\bibitem{chornomaz2025spherical}
B.~Chornomaz, S.~Moran, and T.~Waknine,
\emph{Spherical dimension},
Proceedings of the Thirty Eighth Conference on Learning Theory,
Proc. Mach. Learn. Res. \textbf{291} (2025), 1259--1313.
\href{https://proceedings.mlr.press/v291/chornomaz25a.html}{PMLR 291}.

\bibitem{FLM}
T.~Figiel, J.~Lindenstrauss, and V.~D.~Milman,
\emph{The dimension of almost spherical sections of convex bodies},
Acta Math. \textbf{139} (1977), 53--94.
\href{https://doi.org/10.1007/BF02392234}{doi:10.1007/BF02392234}.

\bibitem{FPP}
P.~Frankl, J.~Pach, and D.~P\'alv\"olgyi,
\emph{Exchange properties of finite set-systems},
SIAM J. Discrete Math. \textbf{36} (2022), 2073--2081.
\href{https://doi.org/10.1137/21M145149X}{doi:10.1137/21M145149X}.

\bibitem{FHV}
F.~Frick, K.~Hosseini, and A.~Vasileuski,
\emph{A $\mathbb Z_2$-topological framework for sign-rank lower bounds},
preprint (2026),
\href{https://arxiv.org/abs/2604.01510}{arXiv:2604.01510}.

\bibitem{GMP}
D.~Govc, W.~Marzantowicz, and P.~Pave\v{s}i\'c,
\emph{Estimates of covering type and the number of vertices
of minimal triangulations},
Discrete Comput. Geom. \textbf{63} (2020), 31--48.
\href{https://doi.org/10.1007/s00454-019-00092-z}{doi:10.1007/s00454-019-00092-z}.

\bibitem{Gromov}
M.~Gromov,
Dimension, non-linear spectra and width,
in \emph{Geometric Aspects of Functional Analysis (1986--87)},
J.~Lindenstrauss and V.~D.~Milman (eds.),
Lecture Notes in Mathematics, vol.~1317,
Springer-Verlag, Berlin, 1988, pp.~132--184.

\bibitem{Hatcher}
A.~Hatcher,
\emph{Algebraic Topology},
Cambridge University Press, 2002.

\bibitem{Haussler1995}
D.~Haussler,
\emph{Sphere packing numbers for subsets of the Boolean $n$-cube
with bounded Vapnik--Chervonenkis dimension},
J. Combin. Theory Ser. A \textbf{69} (1995), 217--232.

\bibitem{HornJohnson}
R.~A.~Horn and C.~R.~Johnson,
\emph{Matrix Analysis},
Cambridge University Press, Cambridge, 1985.

\bibitem{KW}
M.~Karoubi and C.~Weibel,
\emph{On the covering type of a space},
Enseign. Math. \textbf{62} (2016), 457--474.
\href{https://arxiv.org/abs/1612.00532}{arXiv:1612.00532}.

\bibitem{Kim}
J.~H.~Kim,
\emph{The Ramsey number $R(3,t)$ has order of magnitude $t^2/\log t$},
Random Structures Algorithms \textbf{7} (1995), 173--207.
\href{https://doi.org/10.1002/rsa.3240070302}{doi:10.1002/rsa.3240070302}.

\bibitem{Matousek}
J.~Matou\v sek,
\emph{Using the Borsuk--Ulam Theorem: Lectures on Topological Methods
in Combinatorics and Geometry}, Universitext, Springer, 2003.

\bibitem{Milnor}
J.~Milnor, \emph{Morse Theory}, Annals of Mathematics Studies 51,
Princeton University Press, 1963.

\bibitem{Milo}
T.~Milo, \emph{On the Figiel--Lindenstrauss--Milman inequality},
preprint (2025),
\href{https://arxiv.org/abs/2504.13571}{arXiv:2504.13571}.

\bibitem{Sauer}
N.~Sauer,
\emph{On the density of families of sets},
J. Combin. Theory Ser. A \textbf{13} (1972), 145--147.
\href{https://doi.org/10.1016/0097-3165(72)90019-2}{doi:10.1016/0097-3165(72)90019-2}.

\bibitem{ST}
G.~Simonyi and G.~Tardos,
\emph{Local chromatic number, Ky Fan's theorem, and circular colorings},
Combinatorica \textbf{26} (2006), 587--626.
\href{https://doi.org/10.1007/s00493-006-0034-x}{doi:10.1007/s00493-006-0034-x}.

\bibitem{Stanley}
R.~P.~Stanley,
\emph{On the number of faces of centrally-symmetric simplicial polytopes},
Graphs Combin. \textbf{3} (1987), 55--66.
\href{https://doi.org/10.1007/BF01788529}{doi:10.1007/BF01788529}.

\end{thebibliography}
\end{document}
